В этом разделе мы последовательно разбираем основные подходы к эмбеддингу диаграмм персистентности: сначала три классические предписанные конструкции (Persistence Image, Persistence Landscape, Persistence Silhouette), затем два обучаемых подхода (DeepSets и Persformer). Для каждого метода мы кратко описываем устройство, сильные стороны и ограничения с точки зрения сформулированных в §\ref{sec:embedding-task} критериев (i)--(v).

\subsection{Persistence Image}
\label{ssec:pi}

Persistence Image~\cite{adamsPersistenceImages2017} превращает диаграмму $D$ в небольшое <<изображение>> --- двумерную карту значений на регулярной сетке. Сначала диаграмма переводится из координат \emph{birth--death} в координаты \emph{birth--persistence} линейным преобразованием $T(b, d) = (b, d - b)$, после чего по преобразованной диаграмме $T(D)$ строится \emph{поверхность персистентности}:
\[
\rho_D(z) = \sum_{u \in T(D)} f(u) \cdot \phi_u(z),
\]
где $\phi_u$ --- нормированный гауссиан с центром $u$ и фиксированной дисперсией $\sigma^2$, а $f \colon \mathbb{R}^2 \to \mathbb{R}_{\ge 0}$ --- весовая функция, равная нулю на горизонтальной оси (соответствующей диагонали $b = d$) и плавно растущая с увеличением persistence. Persistence image $I(D)$ получается дискретизацией поверхности по равномерной сетке размера $n \times n$ и интегрированием $\rho_D$ по каждой ячейке.

Persistence Image удовлетворяет критериям (i), (ii) и (iii): результат --- вектор фиксированного размера $n^2$, инвариантный по перестановкам точек, а 1-Wasserstein-устойчивость относительно вариаций входа доказана авторами. К ограничениям относятся: необходимость подбирать резолюцию $n$, дисперсию $\sigma^2$ и весовую функцию $f$ под задачу (теряется адаптивность); отсутствие моделирования взаимодействий между точками --- влияние каждой точки <<суммируется>> аддитивно; результат лежит в евклидовом пространстве, но не настраивается под downstream-задачу.

\subsection{Persistence Landscape}
\label{ssec:pl}

Persistence Landscape~\cite{bubenikStatisticalTopologicalData2015} представляет диаграмму как набор кусочно-линейных функций. Для каждой точки $(b_i, d_i)$ диаграммы определяется <<шапочка>>
\[
f_{(b_i, d_i)}(x) = \max\bigl(0, \min(x - b_i, d_i - x)\bigr),
\]
а ландшафт уровня $k$ --- это огибающая $k$-й по убыванию функции:
\[
\lambda_k(x) = \mathop{\mathrm{kmax}}_{i} f_{(b_i, d_i)}(x),
\]
где $\mathop{\mathrm{kmax}}$ обозначает $k$-е по убыванию значение. Ландшафт $\lambda = (\lambda_1, \lambda_2, \dots)$ лежит в $L^p$-Банаховом пространстве, что делает его удобным для статистических процедур: средние, доверительные интервалы и ядра определяются естественно. На практике первые $K$ слоёв оцениваются на равномерной сетке из $R$ точек, что даёт вектор фиксированной длины $K \cdot R$.

Главные преимущества landscape --- наличие сильной теории устойчивости (по $L^\infty$-норме и относительно bottleneck-расстояния) и линейность по диаграммам, позволяющая корректно усреднять. Ограничения схожи с PI: эмбеддинг не имеет обучаемых параметров; число слоёв $K$ и резолюция $R$ выбираются вручную; представление сужается, когда у диаграмм много <<близких>> точек --- они обмениваются местами на разных слоях и могут создавать высокочастотные артефакты.

\subsection{Persistence Silhouette}
\label{ssec:ps}

Persistence Silhouette~\cite{chazalStochasticConvergencePersistence2014} сглаживает диаграмму до единственной функции одной переменной --- взвешенного среднего тех же шапочек $f_{(b_i, d_i)}$:
\[
\phi^{(p)}(x) = \frac{\sum_i |d_i - b_i|^p \cdot f_{(b_i, d_i)}(x)}{\sum_i |d_i - b_i|^p},
\]
где параметр $p \ge 0$ управляет относительным весом высоко- и низкопересистентных точек: при $p \to 0$ все точки учитываются равноправно, при $p \to \infty$ доминируют наиболее значимые. Дискретизация на $R$ узлов сетки даёт вектор фиксированной длины $R$.

Silhouette гораздо компактнее ландшафта (одна функция вместо $K$ слоёв), сохраняет статистическую трактовку и устойчивость, но при этом теряет часть информации о структуре диаграммы --- в частности, не позволяет различить ситуацию с одной значимой точкой от ситуации с несколькими равноценными. Как и landscape, не обучается под задачу.

\subsection{DeepSets}
\label{ssec:deepsets}

DeepSets~\cite{zaheerDeepSets2017} --- общая конструкция для построения перестановочно-инвариантных функций на множествах. Утверждение состоит в том, что любую непрерывную перестановочно-инвариантную функцию на множествах ограниченного размера можно сколь угодно точно приблизить выражением
\[
L(\{x_1, \dots, x_n\}) = \rho \!\left( \sum_{i=1}^{n} \phi(x_i) \right),
\]
где $\phi$ и $\rho$ --- нейронные сети (обычно MLP). В применении к диаграммам персистентности роль $x_i$ играют точки $(b_i, d_i)$ (возможно, с одношаговым one-hot-кодированием гомологической размерности): $\phi$ применяется к каждой точке независимо, результаты симметрично агрегируются (суммой или средним, в реализации часто с маской для padding'а), и итог пропускается через декодирующий MLP $\rho$.

DeepSets обучаема, перестановочно-инвариантна, дифференцируема и масштабируется линейно по числу точек. Главное её ограничение --- отсутствие явного моделирования \emph{взаимодействий} между точками: каждая точка обрабатывается изолированно, а вся <<коммуникация>> сводится к одному pooling-шагу. Для диаграмм, где относительное расположение точек несёт сигнал (например, кластеризация близких к диагонали точек или линейное расположение точек одного и того же масштаба), это становится бутылочным горлышком.

\subsection{Persformer}
\label{ssec:persformer}

Persformer~\cite{reinauerPersformerTransformerArchitecture2021} --- наиболее <<современный>> по архитектуре из рассмотренных подходов. Точки диаграммы трактуются как множество векторов в $\mathbb{R}^2$ (расширенное one-hot-кодированием гомологической размерности), к ним применяется обучаемое позиционное кодирование (линейный input embedding), затем стек self-attention блоков из стандартного трансформера, и наконец --- multi-head attention pooling layer с единственной обучаемой query~\cite{leeSetTransformer2019} для получения вектора фиксированного размера. Полученный вектор обрабатывается небольшим MLP-классификатором.

Главное достоинство Persformer --- self-attention механизм моделирует попарные взаимодействия между всеми точками диаграммы и, как показано авторами, удовлетворяет универсальной аппроксимационной теореме для функций на мультимножествах. Это снимает основное ограничение DeepSets. Кроме того, благодаря saliency maps Persformer допускает интерпретацию: можно оценивать вклад каждой точки диаграммы в финальное предсказание.

Главные ограничения связаны с вычислительной стоимостью: вычислительная сложность self-attention квадратична по числу точек $M$ во входе, $O(M^2)$. Для диаграмм из 30--50 точек это не проблема, однако в нашей задаче с PHT-фильтрацией по 16 направлениям одно изображение порождает сотни и тысячи токенов, и квадратичная сложность становится узким местом и по памяти, и по времени.

\subsection{Сводка}
Таблица~\ref{tab:methods-summary} обобщает сравнение по сформулированным критериям. Видно, что предписанные методы дёшевы и устойчивы, но негибки; DeepSets и Persformer обучаемы и выразительны, но различаются по способности моделировать взаимодействия и по масштабируемости. Именно поиск точки компромисса в треугольнике <<выразительность $\leftrightarrow$ сложность $\leftrightarrow$ масштабируемость>>, с учётом особенностей нашей PHT-постановки, мотивирует архитектуру, описанную в §3.

\begin{table}[ht]
\centering
\fbox{\parbox{0.85\linewidth}{\centering [PLACEHOLDER]\\[0.3em]
Сводная таблица: для каждого метода (PI, PL, PS, DeepSets, Persformer) --- столбцы <<обучаемость>>, <<взаимодействия между точками>>, <<сложность по числу точек $M$>>, <<доказанная устойчивость>>, <<интерпретируемость>>.}}
\caption{Сравнение методов эмбеддинга диаграмм персистентности.}
\label{tab:methods-summary}
\end{table}
